\documentclass[11pt,a4paper]{article}
\usepackage[T1]{fontenc}
\usepackage[utf8]{inputenc}
\usepackage{mathpazo}
\usepackage[scaled=0.92]{helvet}
\usepackage{microtype}
\usepackage[margin=1.1in]{geometry}
\usepackage{amsmath,amssymb,amsthm}
\usepackage{booktabs,tabularx,longtable}
\usepackage{enumitem}
\usepackage[most]{tcolorbox}
\usepackage{titlesec}
\usepackage[hidelinks]{hyperref}
\usepackage{setspace}
\setstretch{1.08}

\titleformat{\section}{\normalfont\Large\bfseries}{\thesection}{1em}{}
\titleformat{\subsection}{\normalfont\large\bfseries}{\thesubsection}{1em}{}
\titleformat{\subsubsection}{\normalfont\normalsize\bfseries\itshape}{\thesubsubsection}{1em}{}

\theoremstyle{plain}
\newtheorem{proposition}{Proposition}[section]
\newtheorem{hypothesis}{Hypothesis}[section]
\theoremstyle{definition}
\newtheorem{definition}{Definition}[section]
\newtheorem{prediction}{Prediction}[section]
\theoremstyle{remark}
\newtheorem{remark}{Remark}[section]

\newtcolorbox{caution}[1][]{colback=gray!6,colframe=gray!55,boxrule=0.5pt,arc=1pt,
  left=6pt,right=6pt,top=4pt,bottom=4pt,fonttitle=\bfseries\small,title={Caution#1}}
\newtcolorbox{statusbox}[1]{colback=white,colframe=black!70,boxrule=0.6pt,arc=0pt,
  left=8pt,right=8pt,top=6pt,bottom=6pt,fonttitle=\bfseries,title={#1}}

\newcommand{\tier}[1]{\textup{\textsc{[#1]}}}
\newcommand{\Eone}{\tier{E1}}
\newcommand{\Etwo}{\tier{E2}}
\newcommand{\Ethree}{\tier{E3}}
\newcommand{\Efour}{\tier{E4}}

\title{\textbf{Mechanical Inheritance}\\[0.4em]
\large Collagen, Prenatal Movement, Maternal Matrix, and the Extracellular Record of Development and Aging\\[0.3em]
\normalsize A Formal Statement of the Research Program, Its Evidentiary Standing, and the Numerous Respects in Which It Remains Unestablished}
\author{Flyxion\\ \small Independent Researcher}
\date{24 September 2026\\ \small Expanded from the research outline of the same date}

\begin{document}
\maketitle

\begin{abstract}
\noindent
This document states, in expanded and deliberately formal terms, a research program connecting four bodies of inquiry: the biology of collagen and the extracellular matrix; Motile Womb Theory, which holds that prenatal movement supplies structured sensorimotor experience and is causally consequential for developing form; the regulation of sexually dimorphic tissue organization; and the role of matrix maintenance in aging and longevity. A fifth, considerably more speculative element, the hypothesis of innate flying dreams, is retained but formally subordinated. The document introduces two organizing proposals not present in the earlier outline. The first is that the mechanical boundary of the prenatal environment is itself a hormonally remodeled maternal collagen structure, so that sexually differentiated matrix regulation enters the theory through the maternal side, where mammalian evidence is comparatively strong. The second is that development and aging may be described along a single axis relating accumulated matrix history to remaining remodeling capacity, such that the property which renders extracellular matrix a durable record also, in time, renders it rigid. Every claim is assigned an explicit evidentiary tier. The document contains no new empirical data, makes no clinical recommendations, and should be read throughout as a map of questions rather than a statement of findings.
\end{abstract}

\begin{statusbox}{Statement of Epistemic Status}
\small
\begin{enumerate}[leftmargin=1.4em,itemsep=2pt]
\item This is a theoretical and programmatic document prepared by an independent researcher without institutional affiliation, laboratory access, or clinical training. It has not undergone peer review.
\item No experiments, observations, surveys, or analyses of original data were conducted in its preparation. All empirical content is drawn from published literature and is summarized rather than re-analyzed.
\item Several references added during expansion (notably those concerning relaxin, cervical remodeling, sex differences in tendon collagen, glucosepane, and long-lived collagen pools) were compiled in the course of drafting and should be verified against the primary literature before this document is circulated or cited. The present author accepts that any such citation may be imprecise in volume, page, or year.
\item Nothing in this document constitutes medical, obstetric, nutritional, or therapeutic advice. In particular, nothing herein recommends any change in behavior during pregnancy, any supplement, or any intervention intended to alter collagen, lifespan, or development.
\item The hypothesis of innate flying dreams is phenomenological and currently lacks a direct evidentiary chain of any kind. Its inclusion does not imply endorsement at the level of the better-supported material, and the better-supported material should not be judged by it.
\item Where the text uses words such as ``memory,'' ``record,'' ``grammar,'' or ``calibration,'' these are to be read in the narrowest material or functional sense compatible with the context unless a stronger reading is explicitly argued for.
\end{enumerate}
\end{statusbox}

\tableofcontents
\newpage
\section{Introduction and Scope}

\subsection{Purpose of the document}

The present document has three purposes. The first is to restate the research program described in the outline of 24 September 2026 with sufficient formality that each constituent claim can be located, qualified, and, where possible, exposed to test. The second is to introduce two organizing proposals, the maternal boundary proposal (Section~\ref{sec:maternal}) and the history--plasticity axis (Section~\ref{sec:axis}), which together supply the program with a structural spine that it previously lacked. The third is to record, in more detail than is customary and perhaps more than is strictly necessary, the conditions under which each claim would fail.

The program is not presented as a theory awaiting confirmation. It is presented as an arrangement of questions, some of which are already largely answered by contemporary developmental biology, some of which are open in a well-defined way, and some of which are open in a way that is not yet well defined. A central aim of the document is to prevent movement of claims between these categories without evidence.

\subsection{Motivating observation}

The motivating observation is modest. Developing organisms move before they are anatomically complete, and the forces produced by that movement are transmitted through tissues whose load-bearing and signaling properties depend heavily on collagen-rich extracellular matrices. It is established that movement participates in constructing skeletal form. It is established that extracellular matrix is not inert packing. The program asks what follows when these two facts are held together and followed forward in time, through sexual differentiation and adult maintenance, into aging.

\subsection{What the program does not claim}

In the interest of clarity, the following claims are explicitly disowned at the outset:

\begin{enumerate}[label=(\roman*),itemsep=2pt]
\item that collagen peptide sequences constitute a general genetic ``grammar'' of body plans or of sexual dimorphism;
\item that ingested collagen, collagen hydrolysate, or marine collagen peptides reach developing or aging tissues as intact signaling entities, or that their consumption has any of the effects discussed herein;
\item that fetuses possess episodic memory of the prenatal environment, or that adults recall it;
\item that flying dreams have been shown to originate in prenatal experience;
\item that any intervention on collagen or matrix has been shown to extend human lifespan;
\item that maternal activity during pregnancy should be modified on the basis of anything written here.
\end{enumerate}

Claim (ii) deserves emphasis. An earlier stratum of speculative material within the author's own archive conflated encoded collagen sequence variation with the dietary behavior of hydrolyzed collagen products. That conflation is rejected here without qualification, in keeping with the separately argued position that an ingested peptide is not a destination and that digestion and redistribution sever any simple relation between what is eaten and what is built.

\subsection{Organization}

Section~\ref{sec:conventions} fixes conventions, including the evidentiary tier system used throughout. Sections~\ref{sec:mwt}--\ref{sec:peptide} restate Motile Womb Theory, the collagen mediator hypothesis, and the taxonomy of collagen-derived objects. Section~\ref{sec:maternal} introduces the maternal boundary proposal. Section~\ref{sec:dimorphism} treats sexual dimorphism. Section~\ref{sec:axis} introduces the history--plasticity axis, and Section~\ref{sec:longevity} treats longevity. Section~\ref{sec:dreams} treats the dream hypothesis, including discriminating predictions and natural experiments. Section~\ref{sec:architecture} proposes a revised architecture. Sections~\ref{sec:program}--\ref{sec:ethics} set out the research program, failure conditions, and ethical constraints. Appendices supply a glossary, an evidence register, and formal notes.

\section{Conventions}\label{sec:conventions}

\subsection{Evidentiary tiers}

Every substantive claim in this document is tagged with one of four tiers. The tiers are ordinal and coarse. They describe the present state of support, not the ultimate truth of the claim, and a claim's tier may be revised in either direction.

\begin{definition}[Evidentiary tiers]
\leavevmode
\begin{description}[leftmargin=2.2em,itemsep=3pt]
\item[\Eone\ Established.] Supported by convergent experimental evidence in at least one relevant organism, with a mechanism identified at least in outline, and not seriously contested in the relevant literature.
\item[\Etwo\ Plausible.] Supported by adjacent evidence, mechanistically coherent, and consistent with \Eone\ material, but lacking direct demonstration in the form stated.
\item[\Ethree\ Open.] Formulable and in principle testable, but without direct support; adjacent evidence may exist but does not discriminate the claim from alternatives.
\item[\Efour\ Phenomenological.] Motivated by features of experience or by analogy; no evidentiary chain currently connects the claim to measurable biology.
\end{description}
\end{definition}

\begin{caution}
Tier assignments are the author's judgment and are not the output of a formal evidence synthesis. No systematic review, risk-of-bias assessment, or meta-analysis has been performed. Readers with domain expertise may reasonably assign different tiers, and should.
\end{caution}

\subsection{Organismal scope}

Claims tagged \Eone\ frequently rest on evidence from \emph{Drosophila melanogaster}, \emph{Caenorhabditis elegans}, zebrafish, chick, or mouse. Transfer of such claims to humans is never assumed. Where a claim is stated for a model organism, it is to be read as holding for that organism unless otherwise specified.

\subsection{Notation}

Throughout, $t$ denotes developmental or chronological time; $u(t)$ denotes a movement signal (fetal, maternal, or combined); $\sigma(x,t)$ denotes local mechanical stress in tissue at position $x$; $m(x,t)$ denotes a vector of matrix-state variables (composition, alignment, crosslink density, turnover rate, ligand presentation); and $c(x,t)$ denotes cellular state. Formal expressions are schematic. They are intended to make dependency structure explicit, not to be fitted to data in their present form.

\begin{caution}
No claim in this document depends on the mathematical expressions being quantitatively correct. The formalism is an aid to stating which variables depend on which, and to identifying where an experiment would need to intervene. It should not be mistaken for a model.
\end{caution}

\section{Motile Womb Theory}\label{sec:mwt}

\subsection{Statement}

Motile Womb Theory proposes that prenatal movement is not merely an output of development but an input to it. Movement supplies structured sensorimotor experience through which the developing organism begins to calibrate orientation, bodily extent, three-dimensional displacement, resistance, acceleration, and the distinction between self-produced and externally produced motion.

\begin{proposition}[Moving environment]\label{prop:moving}
The prenatal organism develops within an environment set in motion by fetal movement, maternal movement, fluid displacement, uterine constraint, gravity, and contact with surrounding tissues. \Eone
\end{proposition}

\begin{proposition}[Correlated input]\label{prop:correlated}
These motions provide correlated vestibular, proprioceptive, tactile, interoceptive, and motor information before postnatal vision and locomotion become dominant. \Eone\ for the existence of such input; \Etwo\ for its correlational structure being developmentally used.
\end{proposition}

\begin{hypothesis}[Primitive calibration]\label{hyp:calibration}
Repeated correlations between action and sensation contribute to primitive spatial calibration and to body--world differentiation. \Etwo
\end{hypothesis}

\begin{hypothesis}[Reuse of early sensorimotor structure]\label{hyp:reuse}
Later spatial imagination and characteristic dream motifs reuse structure established by early sensorimotor calibration. \Efour
\end{hypothesis}

\subsection{Established support}

\paragraph{Movement is morphogenetically consequential.} Reduced or absent fetal movement is associated with abnormal skeletal development, including joint contractures, altered joint shape, and delayed ossification \cite{nowlan2010}. Experimental work identifies mechanosensitive pathways, including PIEZO channels, through which movement-derived forces influence developing bone \cite{zhou2020}. Recent mouse work indicates that mechanical stimulation associated with maternal exercise can partially rescue joint and bone defects produced by fetal akinesia \cite{panebianco2025}. The strong form of the premise, that movement participates causally in constructing the body rather than exercising a completed anatomy, is therefore \Eone\ for skeletal tissues in the organisms studied.

\begin{caution}
The rescue result concerns a mouse model of fetal akinesia under experimental conditions. It does not imply that maternal exercise affects typical human fetal development, nor that exercise can treat any human congenital condition. It is cited solely as evidence that externally supplied mechanical input can reach and influence fetal skeletal tissue in principle.
\end{caution}

\paragraph{The uterus is a constrained mechanical environment.} Biomechanical modeling shows that uterine geometry, fetal position, available space, tissue contact, and amniotic fluid affect the forces generated by fetal movement \cite{verbruggen2016}. The prenatal environment is neither free flight nor weight-bearing locomotion. It is a deformable, fluid-mediated, increasingly constrained field. This supplies a precise reading of ``motile womb'': a time-varying boundary-value problem in which fetal and maternal motion jointly determine mechanical input.

\paragraph{Vestibular development has sensitive periods.} Mammalian vestibular development proceeds through prenatal and postnatal stages, and altered-gravity research supports sensitive-period effects on vestibular structure and function \cite{jamon2014}. This renders prenatal motion a plausible source of early calibration (Hypothesis~\ref{hyp:calibration}).

\begin{caution}
No evidence cited here shows that fetuses construct an adult-like three-dimensional cognitive map, that prenatal vestibular input is necessary for any specific later spatial ability in humans, or that prenatal calibration has any bearing on dreaming. Hypothesis~\ref{hyp:reuse} is addressed separately in Section~\ref{sec:dreams} and is not supported by the material in this subsection.
\end{caution}

\section{The Collagen Mediator Hypothesis}\label{sec:mediator}

\subsection{Statement}

\begin{hypothesis}[Matrix mediation]\label{hyp:mediation}
Developmental movement applies forces to tissues; extracellular matrices distribute, resist, and retain the consequences of those forces; cells sense the resulting matrix state; and remodeling alters the subsequent possibilities of growth and movement. \Etwo
\end{hypothesis}

Schematically,
\[
u(t)\;\to\;\sigma(x,t)\;\to\;\Delta m(x,t)\;\to\;c(x,t)\;\to\;\text{remodeling}\;\to\;\text{altered }u\text{ and growth}.
\]

A minimal dynamical sketch makes the loop explicit:
\begin{align}
\sigma(x,t) &= \mathcal{S}\big[u(t),\, m(x,t),\, g(x,t)\big], \label{eq:stress}\\
\partial_t m(x,t) &= \mathcal{R}\big[m(x,t),\, \sigma(x,t),\, c(x,t)\big], \label{eq:remodel}\\
\partial_t c(x,t) &= \mathcal{C}\big[c(x,t),\, m(x,t),\, \sigma(x,t),\, \phi(x,t)\big], \label{eq:cell}
\end{align}
where $g$ denotes tissue geometry and $\phi$ denotes morphogen and endocrine fields. The feature of interest is that $m$ appears on both sides: matrix state conditions the stress it experiences, and stress conditions how matrix state changes.

\begin{remark}
The formulation is deliberately stronger than a claim that collagen sequence encodes form, and weaker than a claim that collagen is the principal mediator of mechanical effects in development. Other matrix components (proteoglycans, elastin, fibronectin, laminins), cytoskeletal mechanics, and cell--cell junctions carry mechanical information as well. Collagen is foregrounded because of its abundance and load-bearing role, not because it is presumed to act alone.
\end{remark}

\subsection{Established support}

Extracellular matrix provides scaffolding, force transmission, growth-factor sequestration, receptor engagement, and spatial constraint on migration and deformation. Collagen production and processing during embryogenesis are tissue-specific. Homologous ADAMTS proteases acquire distinct developmental roles through spatial correspondence with particular procollagen substrates \cite{legoff2006}. Collagen XV affects notochord differentiation and muscle development in zebrafish and interacts with Hedgehog-dependent patterning \cite{pagnon2008}. Embryo attachment and outgrowth \emph{in vitro} depend on collagen type and on exposure of binding sequences such as Arg-Gly-Asp \cite{armant1988}. The conclusion that collagen matrices form part of the causal environment of development is \Eone.

\begin{caution}
That matrix is causally involved in development does not establish that matrix mediates the specific long-term consequences of \emph{movement}. Hypothesis~\ref{hyp:mediation} remains \Etwo\ until movement manipulation is shown to alter matrix state and until that altered matrix state is shown to carry forward the downstream morphological effect. Demonstrating the first without the second would leave open that movement acts through cellular pathways that bypass matrix entirely.
\end{caution}

\section{The Peptide Question}\label{sec:peptide}

\subsection{Disambiguation}

The phrase ``collagen peptide'' is ambiguous to a degree that has, in earlier stages of this program, produced genuine confusion. It is resolved here into five categories, of which only the first four are biologically relevant to the hypotheses of this document.

\begin{center}
\small
\begin{tabularx}{\textwidth}{@{}>{\raggedright\arraybackslash}p{0.27\textwidth}>{\raggedright\arraybackslash}X>{\raggedright\arraybackslash}p{0.2\textwidth}@{}}
\toprule
\textbf{Category} & \textbf{Description} & \textbf{Relevance} \\
\midrule
C1. Encoded chains and domains & Gene products whose sequences govern folding, trimer selection, receptor binding, fibrillogenesis, and tissue specialization & Central \\
C2. Propeptides and processing products & Cleaved during biosynthesis; possible regulators and established biomarker candidates & Central \\
C3. Endogenous fragments (matrikines) & Released by proteolysis during remodeling; some bioactive \cite{maquart2005} & Plausible, underdeveloped \\
C4. Synthetic binding or hybridizing peptides & Laboratory and therapeutic tools that target damaged or unfolded collagen & Methodological \\
C5. Dietary hydrolysates & Commercial products; digested prior to absorption & \emph{Excluded} \\
\bottomrule
\end{tabularx}
\end{center}

\begin{caution}
Category C5 is listed only in order to exclude it. Nothing in this program depends on, supports, or should be cited in support of claims regarding collagen supplements. Any reader encountering this document in connection with such products should treat that connection as a misreading.
\end{caution}

\subsection{Requirements for a sequence-level program}

A sequence-level developmental claim concerning C3 would require, at minimum: identification of the exact peptide; identification of the protease responsible for its release; measurement of its local concentration; identification of its receptor; demonstration of receptor expression in the relevant tissue at the relevant time; and a perturbation showing that removal or blockade of the peptide alters the developmental outcome. None of these requirements is currently met for any claim advanced by this program. Sequence-level claims are therefore uniformly \Ethree.
\section{The Maternal Boundary Proposal}\label{sec:maternal}

\subsection{Motivation}

Section~\ref{sec:mwt} described the prenatal environment as a boundary-value problem. The earlier outline treated the boundary as given. It is not given. The uterine wall, the cervix, the pelvic ligaments, and the connective tissues surrounding them are themselves collagen-rich matrices, and during pregnancy they undergo extensive and hormonally directed remodeling. The mechanical boundary conditions of the motile womb are therefore partly determined by a maternal matrix whose properties change across gestation.

\subsection{Statement}

\begin{proposition}[Maternal matrix remodeling in pregnancy]\label{prop:maternalremodel}
Pregnancy involves substantial remodeling of maternal reproductive and pelvic connective tissue, including progressive softening and reorganization of cervical collagen \cite{word2007}, under the influence of hormones including relaxin and estrogens that act on collagen synthesis, degradation, and organization \cite{sherwood2004}. \Eone\ for the existence of this remodeling in mammals; the details vary considerably across species.
\end{proposition}

\begin{hypothesis}[Maternal boundary]\label{hyp:maternalboundary}
The mechanical input reaching the fetus depends in part on the time-varying compliance and geometry of maternal matrix, such that maternal matrix state functions as a boundary condition on fetal movement and fetal tissue loading. \Etwo
\end{hypothesis}

Formally, let $\Omega(t)$ denote the intrauterine domain and $\partial\Omega(t)$ its boundary. Fetal tissue stress $\sigma_f$ satisfies a mechanical problem of the schematic form
\[
\nabla\!\cdot\sigma_f = \mathbf{f}(u_f, u_m) \quad\text{in } \Omega(t), \qquad
\sigma_f\cdot\mathbf{n} = \mathcal{B}\big[m_{\mathrm{mat}}(t)\big]\quad\text{on }\partial\Omega(t),
\]
where $u_f$ and $u_m$ are fetal and maternal movement and $\mathcal{B}$ maps maternal matrix state $m_{\mathrm{mat}}$ to an effective boundary compliance. The substantive content of Hypothesis~\ref{hyp:maternalboundary} is that $\mathcal{B}$ is not constant and that its variation is non-negligible for fetal loading.

\begin{caution}
The expression above is illustrative only. Real intrauterine mechanics involve fluid--structure interaction, placental attachment, fetal position changes, and heterogeneous maternal tissues, none of which are represented. No quantitative claim is made about the magnitude of maternal matrix effects on fetal loading, which may prove small relative to amniotic fluid volume, fetal size, or position.
\end{caution}

\subsection{Consequence for the architecture of the program}

The maternal boundary proposal alters the route by which sexual dimorphism enters the program. The earlier outline sought a connection between collagen and sexual dimorphism primarily through the developing organism's own sex-specific matrix regulation, where mammalian evidence is thin. Under Hypothesis~\ref{hyp:maternalboundary}, a second route exists: a sex-specific, hormonally directed program of matrix remodeling in the \emph{mother} shapes the mechanical environment of \emph{every} fetus, regardless of fetal sex. The phenomenon connecting sexually differentiated matrix biology to the motile womb is thereby moved from \Ethree\ territory into a region adjacent to \Eone\ material.

\begin{remark}
This route does not explain sexual dimorphism in offspring. It explains how sexually differentiated matrix biology enters the developmental environment. The two questions must not be merged. A maternal program that shapes all fetuses equally contributes to development, not to differences between the sexes of offspring.
\end{remark}

\subsection{A threefold coupling}

Held together, Propositions~\ref{prop:maternalremodel} and Hypotheses~\ref{hyp:mediation} and~\ref{hyp:maternalboundary} describe a coupling of three matrices across two organisms:
\begin{align*}
&\underbrace{\text{maternal hormones}}_{\text{endocrine}}
\;\to\;
\underbrace{m_{\mathrm{mat}}(t)}_{\text{maternal matrix}}
\;\to\;
\underbrace{\partial\Omega(t)}_{\text{boundary}}\\
&\qquad\to\;
\underbrace{\sigma_f}_{\text{fetal loading}}
\;\to\;
\underbrace{m_f(t)}_{\text{fetal matrix}}.
\end{align*}
The fetal matrix is thus shaped by forces constrained by another organism's matrix, which is itself shaped by a hormonal program. This is the single most unifying structure available to the program at present.

\begin{caution}
Establishing this chain in humans would face severe practical and ethical limits (Section~\ref{sec:ethics}). The chain is proposed as an organizing hypothesis for animal and modeling work, not as a claim about human pregnancy.
\end{caution}

\section{Sexual Dimorphism}\label{sec:dimorphism}

\subsection{Invertebrate precedent}

In \emph{Drosophila}, sex-specific isoforms of the transcription factor Doublesex regulate the production and secretion of collagen IV, partly through procollagen lysyl hydroxylase, allowing a sex-determination pathway to alter basement membranes beyond the cells in which the initiating factor is expressed \cite{peng2022}.

\begin{proposition}[Sex-determination control of collagen]\label{prop:dsx}
In at least one well-characterized animal, sex-determination machinery regulates collagen production and secretion and thereby influences tissue environments non-cell-autonomously. \Eone\ for \emph{Drosophila}.
\end{proposition}

\begin{caution}
Mammalian sex determination is mechanistically distinct from that of \emph{Drosophila}. Doublesex-related (DMRT) genes exist in mammals, but no inference is drawn here that they regulate collagen in the same manner. Proposition~\ref{prop:dsx} is cited as an existence proof of a category of mechanism, not as evidence for any mammalian instance.
\end{caution}

\subsection{Mammalian candidate systems}

Three mammalian systems are proposed as tractable starting points, each chosen because sex-differentiated or hormonally modulated collagen biology is already documented there in some form.

\paragraph{Tendon and ligament.} Sex and estrogen appear to influence musculotendinous collagen synthesis at rest and following exercise \cite{hansen2014}. This supplies a system in which hormonally modulated matrix turnover can be studied alongside mechanical loading, which is precisely the conjunction the program requires. \Etwo\ as a model for the program's purposes.

\paragraph{Pelvic and reproductive connective tissue.} The remodeling described in Proposition~\ref{prop:maternalremodel} is the most dramatic sex-limited matrix program in mammals. \Eone\ for existence.

\paragraph{Dermis and skeletal tissues.} Sex differences in skin and bone properties are widely reported, but their attribution to collagen regulation specifically, rather than to hormonal effects on cells, body size, or loading history, is not straightforward. \Ethree\ as a collagen-mediated effect.

\subsection{The causal chain required}

\begin{hypothesis}[Matrix-mediated dimorphism]\label{hyp:dimorphism}
Sexually differentiated regulators or hormones alter collagen synthesis or processing, thereby altering matrix properties, thereby contributing to a dimorphic tissue outcome:
\begin{align*}
&\text{sex regulator or hormone}\;\to\;\text{collagen synthesis or processing}\\
&\qquad\to\;\text{matrix property}\;\to\;\text{dimorphic outcome}.
\end{align*}
\Etwo\ for the chain's existence in some tissue; \Ethree\ for its general importance to body form.
\end{hypothesis}

\begin{caution}
Comparative genomics alone cannot establish Hypothesis~\ref{hyp:dimorphism}. Correlation between collagen sequence variation and degree of dimorphism across species (for instance, across taxa with extreme dimorphism such as anglerfish) would be confounded by phylogeny, body size, ecology, and life history. Establishing each arrow requires perturbation, rescue, tissue-specific manipulation, and attention to developmental timing.
\end{caution}

\subsection{Explicitly retired claim}

The earlier speculative claim that ``collagen peptide sequences systematically organize sexual dimorphism across species'' is retained in the evidence register only as an \Ethree\ item and is not used as a premise anywhere in this document. Its original formulation did not distinguish categories C1 through C5 (Section~\ref{sec:peptide}) and should not be cited in its original form.

\section{The History--Plasticity Axis}\label{sec:axis}

\subsection{Motivation}

The earlier outline proposed that development and aging occupy opposite portions of a continuous matrix-maintenance problem. This section attempts to make that proposal precise by identifying a single axis along which both can be located.

\subsection{Two quantities}

\begin{definition}[Recorded matrix history]
For a tissue region $X$ at time $t$, the recorded history $H_X(t)$ is the extent to which the present matrix state of $X$ reflects past loading, damage, and chemical modification rather than recent synthesis. Candidate proxies include the fraction of collagen older than a threshold age, advanced glycation end-product crosslink density, and residual alignment attributable to past loading.
\end{definition}

\begin{definition}[Remodeling capacity]
The remodeling capacity $P_X(t)$ is the rate at which the matrix of $X$ can be replaced, reorganized, or repaired in response to present conditions. Candidate proxies include collagen turnover rate, protease and synthetic enzyme activity, and responsiveness of matrix synthesis to mechanical loading.
\end{definition}

\begin{definition}[History--plasticity ratio]
\[
\Lambda_X(t) \;=\; \frac{H_X(t)}{P_X(t)}.
\]
\end{definition}

\begin{remark}
$\Lambda$ is not proposed as a measurable scalar in its present form. $H$ and $P$ are each multidimensional, and their proxies are not commensurable without further work. The ratio is introduced to name a direction, not to supply a number.
\end{remark}

\subsection{Supporting observations}

Human matrix contains pools that turn over extremely slowly. Collagen in articular cartilage has been estimated to have a half-life on the order of a century \cite{verzijl2000}, and radiocarbon evidence indicates that the core of the adult human Achilles tendon undergoes little renewal after adolescence \cite{heinemeier2013}. Such tissues are, in a literal sense, records: their present composition reflects conditions during their formation. Glucosepane has been identified as a major crosslink in senescent human extracellular matrix \cite{sell2005}, and crosslink accumulation of this kind is associated with increased stiffness. \Eone\ for the existence of long-lived collagen pools and age-associated crosslink accumulation.

\subsection{Statement}

\begin{hypothesis}[Single axis]\label{hyp:axis}
Early development is characterized by low $\Lambda$ (little accumulated history, high remodeling capacity); adult maintenance by intermediate and approximately stable $\Lambda$; and aging by rising $\Lambda$, as history accumulates and remodeling capacity declines. \Etwo\ as a qualitative description; \Ethree\ as an organizing principle with predictive force.
\end{hypothesis}

\begin{proposition}[Record and rigidity]\label{prop:tradeoff}
The same material property, resistance to replacement, that makes extracellular matrix a durable record of past conditions also makes it progressively unable to adapt to present conditions. \Etwo
\end{proposition}

Proposition~\ref{prop:tradeoff} carries the main conceptual weight of the axis. It implies that the persistence of developmental organization and the rigidity of aged tissue are not two phenomena requiring two explanations, but two aspects of a single property viewed at different times. The matrix-as-memory language of the earlier outline (``matrix as developmental memory'') is thereby given a cost: memory in this material sense is purchased with plasticity.

\begin{prediction}\label{pred:axis1}
Tissues with slower collagen turnover should retain signatures of developmental loading history into adulthood more strongly than tissues with rapid turnover, and should also show earlier or steeper age-associated loss of mechanical adaptability. \Ethree
\end{prediction}

\begin{prediction}\label{pred:axis2}
Interventions that extend healthspan in model organisms should, where they act through matrix, tend to preserve remodeling capacity $P$ rather than merely increase collagen quantity. \Ethree
\end{prediction}

\begin{caution}
Hypothesis~\ref{hyp:axis} may be false in informative ways. Some tissues may remain plastic late; some developmental matrices may be highly persistent from the outset; and $\Lambda$ may fail to rise monotonically. Nothing here implies that reducing $H$ (for example, by removing crosslinks) would be beneficial, since accumulated structure also carries function. No therapeutic inference should be drawn.
\end{caution}

\section{Longevity}\label{sec:longevity}

\subsection{Matrix as an aging system}

Aging alters collagen abundance, crosslinking, proteolysis, glycation, organization, and coupling to cells. These changes bear on mechanical feedback, barrier function, repair, stem-cell niches, vascular behavior, and cell--matrix communication. \Eone\ for the existence of these changes; their causal contribution to organismal aging, as distinct from their being consequences of it, is \Etwo.

\subsection{Causal evidence in \emph{C.\ elegans}}

In \emph{C.\ elegans}, several longevity interventions preserve subsets of cuticular collagens and collagen-remodeling enzymes, and disruption of collagen stabilization or matrix coupling can blunt longevity produced by reduced insulin/IGF-1 signaling \cite{ewald2015,teuscher2024}. Experimental collagen overexpression has been reported to extend lifespan in some contexts.

\begin{proposition}[Matrix maintenance in nematode longevity]
Extracellular matrix maintenance is causally involved in some longevity pathways in \emph{C.\ elegans}. \Eone\ for that organism.
\end{proposition}

\begin{caution}
The nematode cuticle is an exoskeletal collagenous structure that differs profoundly from mammalian basement membranes and interstitial matrix. Nothing in the nematode literature establishes that increasing human collagen, consuming collagen, or manipulating human matrix will extend human life. The relevant open question is which principles, if any, are conserved.
\end{caution}

\subsection{Reformulation under the axis}

Under Hypothesis~\ref{hyp:axis}, the longevity question becomes: which interventions hold $\Lambda$ approximately stable for longer? This reformulation shifts attention from matrix quantity to matrix responsiveness, and aligns with the nematode observation that preserved remodeling machinery, not merely preserved collagen, accompanies extended lifespan.

\begin{hypothesis}[Maintenance of coupling]
Longevity requires continued maintenance of mechanical and extracellular coupling between cells and their surroundings, not merely protection of intracellular components. \Etwo\ in nematodes; \Ethree\ in mammals.
\end{hypothesis}

\subsection{Collagen fragments as readouts}

Collagen fragments (C2, C3) may serve as indicators of tissue-specific turnover, and synthetic collagen-targeting peptides (C4) may deliver imaging or therapeutic agents to damaged matrix. Their value depends on distinguishing causal signals from degradation products and on declining to assume that any collagen-derived fragment is beneficial.
\section{Innate Flying Dreams}\label{sec:dreams}

\subsection{Status}

The hypothesis interprets flying dreams as possible reconstructions of prenatal spatial experience: buoyancy without walking, rotation without a stable horizon, whole-body translation generated in part by another organism, and motion within compliant boundaries. It is an original phenomenological proposal and not an established result in dream science. It is assigned \Efour\ in its present form.

\begin{statusbox}{Subordination Notice}
\small
In the architecture proposed in Section~\ref{sec:architecture}, the dream hypothesis is a speculative downstream application of the program and not one of its pillars. Failure of the dream hypothesis would have no bearing on the matrix mediation, maternal boundary, dimorphism, or axis hypotheses. Conversely, support for those hypotheses would provide no support for the dream hypothesis. The two must not be used to validate one another.
\end{statusbox}

\subsection{The admissible bridge}

The only bridge considered admissible is persistence of sensorimotor priors, not persistence of memory.

\begin{hypothesis}[Prior persistence]\label{hyp:prior}
Early vestibular and proprioceptive calibration shapes durable expectations regarding acceleration, orientation, bodily suspension, and controllability, and these expectations are available to generative processes during dreaming, without any narratively accessible memory of the prenatal period being preserved. \Ethree
\end{hypothesis}

\begin{caution}
Hypothesis~\ref{hyp:prior} does not claim that dream content is a replay of anything. It claims only that the space of dream kinematics may be constrained by early calibration. Even if true, it would be consistent with flying dreams having many other determinants, including waking experience, culture, media, and ongoing vestibular activity during sleep.
\end{caution}

\subsection{Discriminating phenomenology}

The principal difficulty identified in the earlier outline is that postnatal experiences (being carried, rocked, swimming, falling, swinging, viewing flight, immersive games) could account for the same dream content. A partial remedy is to specify, in advance, which phenomenological features would be characteristic of a prenatal source and which of a culturally learned one.

\begin{center}
\small
\begin{tabularx}{\textwidth}{@{}XX@{}}
\toprule
\textbf{Features predicted by a prenatal-prior source} & \textbf{Features predicted by a cultural-learning source} \\
\midrule
Passive translation, motion imposed rather than chosen & Directed heading and deliberate navigation \\
Rotation without a stable horizon & Stable horizon, landscape viewed from altitude \\
Compliant enclosure, soft or pressing boundaries & Open sky, unbounded space \\
Buoyancy, suspension, weightlessness without propulsion & Effortful or mechanically explained propulsion (arms, capes, vehicles) \\
Low visual content relative to bodily sensation & Strongly visual, cinematic framing \\
\bottomrule
\end{tabularx}
\end{center}

\begin{prediction}\label{pred:dream1}
If Hypothesis~\ref{hyp:prior} holds, flying and floating dream reports coded blind to hypothesis should cluster on the left column more than base rates of the relevant waking experiences would predict. \Ethree
\end{prediction}

\begin{caution}
The two columns are not exhaustive or exclusive, and dream reports are notoriously shaped by the act of reporting. A clustering on the left column would be weak evidence at best, since compliant enclosure and passive motion are also features of early postnatal care and of vestibular sensation during sleep generally. A clustering on the right column would count against the hypothesis more decisively than the converse would count for it.
\end{caution}

\subsection{Natural experiments}

Because prenatal movement cannot ethically be manipulated in humans, any human test must rely on conditions that vary it naturally. Candidate conditions include:

\begin{enumerate}[itemsep=3pt]
\item \textbf{Reduced amniotic fluid}, which is associated with restricted fetal movement and, in severe cases, with contractures.
\item \textbf{Multiple pregnancy}, in which available intrauterine space and contact conditions differ substantially from singleton pregnancy.
\item \textbf{Persistent breech presentation}, which in some cases may be both a consequence and a marker of altered fetal movement.
\item \textbf{Preterm birth}, which removes the fluid-supported environment weeks before term and substitutes an incubator environment with very different vestibular and tactile properties.
\item \textbf{Congenital vestibular impairment}, which tests whether dream flight and floating depend on vestibular calibration at all, independent of prenatal movement.
\end{enumerate}

\begin{caution}
Every one of these conditions is heavily confounded. Reduced amniotic fluid and breech presentation may reflect underlying fetal or placental conditions; multiple and preterm births are associated with differences in neurological development, early medical care, and family circumstance; congenital vestibular impairment is frequently accompanied by hearing loss and developmental differences that could independently affect dream content or dream reporting. Any observed association would require careful adjustment and would remain correlational. These conditions are listed as possible sources of variation, not as validated instruments. Nothing here characterizes any clinical group, and individuals with these histories should not infer anything about themselves from this section.
\end{caution}

\subsection{Present standing}

The defensible statement remains: prenatal sensorimotor development offers a possible source of constraints on later dream kinematics. It is not a demonstrated origin of flying dreams, and in its present form it is not even well tested.

\section{A Revised Architecture}\label{sec:architecture}

\subsection{The problem with the earlier arrangement}

In the earlier outline, four branches (Motile Womb Theory, collagen, dimorphism, longevity) and one speculative extension (dreams) were joined chiefly by recurrence of the words \emph{collagen} and \emph{movement}. Such lexical coherence is not structural coherence. It also exposes the stronger material to guilt by association with the weaker.

\subsection{Proposed spine}

The program is reorganized around a single material thesis:

\begin{statusbox}{Central Thesis}
Extracellular matrix is the medium through which mechanical history persists in tissue. Its persistence is constructive in development, constraining in maturity, and limiting in age.
\end{statusbox}

Around this thesis, the components are arranged as follows:

\begin{center}
\small
\begin{tabularx}{\textwidth}{@{}>{\raggedright\arraybackslash}p{0.27\textwidth}>{\raggedright\arraybackslash}X>{\raggedright\arraybackslash}p{0.2\textwidth}@{}}
\toprule
\textbf{Component} & \textbf{Role in architecture} & \textbf{Tier} \\
\midrule
Movement morphogenesis & Source of mechanical history & \Eone \\
Matrix mediation (Hyp.~\ref{hyp:mediation}) & Mechanism of persistence & \Etwo \\
Maternal boundary (Hyp.~\ref{hyp:maternalboundary}) & Boundary conditions of the earliest history; entry point for sexually differentiated matrix biology & \Etwo \\
Matrix-mediated dimorphism (Hyp.~\ref{hyp:dimorphism}) & Differentiated regulation of persistence & \Etwo/\Ethree \\
History--plasticity axis (Hyp.~\ref{hyp:axis}) & Continuity from development to aging & \Etwo/\Ethree \\
Longevity & Late portion of the axis & \Eone\ (nematode) / \Ethree\ (mammal) \\
Flying dreams (Hyp.~\ref{hyp:prior}) & Speculative downstream application & \Efour \\
\bottomrule
\end{tabularx}
\end{center}

\begin{remark}
Under this arrangement the program's scientific value does not depend on any \Ethree\ or \Efour\ item. If every such item were refuted, the remaining structure would still constitute a coherent synthesis of existing developmental and aging biology, albeit a less original one.
\end{remark}

\section{Research Program}\label{sec:program}

The following work packages are stated in order of dependency. They are not proposals for which funding or facilities exist. They describe what a competent research group would need to do.

\subsection{WP1. Controlled vocabulary}

Construct a vocabulary distinguishing collagen genes, chains, domains, propeptides, endogenous fragments, post-translational modifications, receptors, proteases, and matrix architectures, keyed to categories C1 through C5. Require every claim in the program to name its object in this vocabulary. This work package can be undertaken without laboratory access and is the only one fully within reach of the present author.

\subsection{WP2. Developmental matrix atlas}

For selected model organisms, map collagen and collagen-processing gene expression across developmental time, sex, tissue, and mechanical context, integrating existing spatial transcriptomic and proteomic resources with morphogen pathway data. Much of the required data may already exist in public repositories.

\subsection{WP3. Movement--matrix coupling}

Using fetal-akinesia, altered-loading, embryo-culture, or organoid systems, determine whether altered movement changes collagen alignment, processing, crosslinking, turnover, or receptor signaling, and whether those changes carry forward to later tissue form. Distinguish mechanical effects from placental, endocrine, nutritional, and maternal-stress effects.

\subsection{WP4. Maternal boundary}

In animal models, determine whether manipulation of maternal matrix compliance (for example, via relaxin signaling) measurably alters fetal loading or fetal skeletal outcomes, independent of hormonal effects acting directly on the fetus. Separating these two routes is the principal technical obstacle.

\subsection{WP5. Dimorphism mechanisms}

Test explicit causal chains of the form given in Hypothesis~\ref{hyp:dimorphism} in tendon, ligament, and reproductive tissue, using perturbation and rescue.

\subsection{WP6. Axis estimation}

Assemble existing data on collagen turnover, crosslink density, and mechanical adaptability across tissues and ages to test whether any consistent ordering corresponding to $\Lambda$ emerges (Predictions~\ref{pred:axis1}--\ref{pred:axis2}).

\subsection{WP7. Dream phenomenology}

Develop and preregister a phenomenological coding scheme for flying and floating dreams (trajectory, controllability, gravity, enclosure, rotation, viewpoint, effort, affect). Only subsequently examine associations with prenatal or early-life movement proxies. This work package proceeds independently of all others.

\section{Failure Conditions}\label{sec:failure}

In the interest of making the program vulnerable, the following observations would count against its components. The list is not exhaustive.

\begin{enumerate}[itemsep=4pt]
\item \textbf{Matrix mediation fails} if movement manipulation alters later morphology while matrix state is shown to be unaltered, or if restoring matrix state does not restore the downstream outcome.
\item \textbf{The maternal boundary proposal fails} if maternal matrix compliance, varied within physiological range, produces no measurable change in fetal loading relative to fluid volume, position, and fetal size.
\item \textbf{Matrix-mediated dimorphism fails}, in a given tissue, if blocking the sex-differentiated matrix response leaves the dimorphic outcome intact.
\item \textbf{The history--plasticity axis fails} if persistence of developmental signatures and loss of adult adaptability are uncorrelated across tissues, or if $\Lambda$ proxies show no consistent age trend.
\item \textbf{The longevity extension fails} if matrix-preserving mechanisms are found to be specific to the nematode cuticle and absent from conserved mammalian matrix biology.
\item \textbf{The dream hypothesis fails} if blind-coded flying dreams cluster on cultural-learning features, or if dream kinematics show no association with any prenatal proxy after adequate adjustment, or if individuals with congenital vestibular impairment report flying and floating dreams indistinguishable from others.
\end{enumerate}

\begin{caution}
The final failure condition is stated with particular diffidence. Dream reports in any population depend on recall, vocabulary, and expectation, and absence of a difference would not be conclusive. It is nevertheless recorded, because a hypothesis for which no failure condition can be stated is not yet a hypothesis.
\end{caution}

\section{Ethical Constraints and Limitations}\label{sec:ethics}

\subsection{Ethical constraints}

Prenatal movement, maternal activity, and pregnancy physiology cannot be manipulated in humans for research purposes of the kind described here. Human evidence must come from observation, natural variation, and existing clinical records, subject to appropriate consent and oversight. Research involving pregnant individuals, preterm infants, or people with congenital conditions carries particular obligations of care, and none of the populations mentioned in Section~\ref{sec:dreams} should be characterized, stigmatized, or burdened on the basis of speculative hypotheses.

This document must not be read as suggesting that pregnant individuals should increase, decrease, or otherwise modify physical activity, nor that fetal outcomes are attributable to maternal behavior. Guidance on activity during pregnancy is a matter for clinical practice and is outside the competence of this document and its author.

\subsection{Limitations}

\begin{enumerate}[itemsep=2pt]
\item The author lacks laboratory, clinical, and obstetric training.
\item No systematic search of the literature was performed; relevant contrary evidence may exist and may be decisive.
\item Tier assignments are judgments, not measurements.
\item The formal expressions are schematic and have not been checked against any mechanical or biochemical data.
\item Several of the cited works were added from recollection during drafting and require verification.
\item The program spans fields in which the author has uneven depth, and the synthesis may overstate the ease with which findings in one field bear on another.
\end{enumerate}

\section{Conclusion}

The program has a coherent core. Development occurs in a mechanically active environment; extracellular matrix can mediate the conversion of movement into persistent biological organization; the boundary of the earliest such environment is itself a hormonally remodeled maternal matrix; and the persistence that makes matrix a record also, over time, makes it rigid. Contemporary biology supports the importance of fetal movement, mechanotransduction, tissue-specific collagen processing, maternal connective-tissue remodeling, long-lived collagen pools, and matrix maintenance in nematode longevity.

The program does not establish that collagen sequences constitute a general grammar of sexual dimorphism, that prenatal motion is encoded in adult dreams, or that any collagen-directed intervention will extend human life. Those are separate hypotheses requiring separate evidence, and they are held at arm's length from the core. The immediate task is to make the middle of the theory explicit and testable: movement changes matrix state; maternal matrix bounds that movement; matrix state changes cellular decisions; and the accumulated record of those decisions both sustains and eventually limits the organism.

\appendix

\section{Glossary}

\begin{description}[leftmargin=1.5em,itemsep=2pt]
\item[Basement membrane.] A thin, specialized extracellular matrix layer, rich in collagen IV and laminins, underlying epithelia and surrounding certain cell types.
\item[Fetal akinesia.] Absence or severe reduction of fetal movement.
\item[Glucosepane.] An advanced glycation end-product crosslink formed between lysine and arginine residues, abundant in aged human matrix.
\item[Matrikine.] A bioactive peptide released by partial proteolysis of extracellular matrix proteins.
\item[Mechanotransduction.] Conversion of mechanical stimuli into biochemical signals by cells.
\item[Procollagen.] The precursor form of fibrillar collagen, bearing propeptides removed during maturation.
\item[Relaxin.] A peptide hormone of pregnancy with documented effects on connective-tissue remodeling in several mammals.
\item[Remodeling capacity ($P$).] As defined in Section~\ref{sec:axis}.
\item[Recorded history ($H$).] As defined in Section~\ref{sec:axis}.
\end{description}

\section{Evidence Register}

\begin{center}
\small
\begin{longtable}{@{}p{0.72\textwidth}l@{}}
\toprule
\textbf{Claim} & \textbf{Tier} \\
\midrule
\endhead
Fetal movement generates forces required for normal skeletal morphogenesis & \Eone \\
Vestibular development is sensitive to environmental input and altered gravity & \Eone \\
Collagen matrices transmit force, organize adhesion, and undergo tissue-specific processing in embryogenesis & \Eone \\
Sex-determination machinery regulates collagen in \emph{Drosophila} & \Eone \\
Pregnancy involves hormonally directed maternal connective-tissue remodeling & \Eone \\
Long-lived collagen pools exist in adult human tissue & \Eone \\
Age-associated crosslinks accumulate in human matrix & \Eone \\
Matrix maintenance is causally involved in some \emph{C.\ elegans} longevity pathways & \Eone \\
Correlated prenatal sensorimotor input is developmentally used & \Etwo \\
Matrix mediates long-term anatomical consequences of prenatal movement & \Etwo \\
Maternal matrix functions as a boundary condition on fetal loading & \Etwo \\
Sex-biased matrix regulation contributes to dimorphic outcomes in some tissues & \Etwo \\
Record and rigidity are aspects of one matrix property & \Etwo \\
Development and aging lie on a single history--plasticity axis & \Etwo/\Ethree \\
Sex-biased matrix regulation broadly shapes dimorphic body plans & \Ethree \\
Specific collagen peptide sequences organize dimorphism across species & \Ethree \\
Endogenous collagen fragments coordinate developmental patterning as proposed & \Ethree \\
Matrix-directed interventions extend mammalian lifespan & \Ethree \\
Prenatal calibration constrains later dream kinematics & \Ethree \\
Flying dreams originate in prenatal spatial experience & \Efour \\
\bottomrule
\end{longtable}
\end{center}

\section{Formal Notes}

\subsection{On the non-commutativity of the mediation loop}

Equations~\eqref{eq:stress}--\eqref{eq:cell} imply that the order of mechanical events matters. If two loading histories $u_1$ and $u_2$ have the same time-integrated magnitude but differ in sequence, the resulting matrix states will in general differ, because each increment of stress acts on a matrix already altered by prior increments. This is the formal content of the claim that matrix records \emph{history} rather than \emph{total exposure}, and it supplies one experimental signature: sequence-permuted loading protocols of equal total magnitude should yield distinguishable matrix outcomes if Hypothesis~\ref{hyp:mediation} holds.

\subsection{On the monotonicity of $\Lambda$}

Suppose $H$ increases whenever the replacement rate falls below the rate of damage and modification, and $P$ declines with age. Then $\Lambda$ rises monotonically once replacement falls below damage. The interesting cases are those in which $\Lambda$ does not rise, which would indicate tissues whose maintenance keeps pace with modification. Identifying such tissues, if they exist, would be as informative for longevity research as confirming the general trend.

\subsection{On the independence of the dream component}

Let $D$ denote the dream hypothesis and $M$ the conjunction of the matrix hypotheses. The architecture of Section~\ref{sec:architecture} is designed so that $P(M \mid \text{evidence})$ is invariant to evidence bearing only on $D$, and vice versa. Any argument that appears to transfer support between them should be treated as a defect in the argument.

\begin{thebibliography}{99}
\small
\bibitem{armant1988} Armant DR, et al. Collagens support embryo attachment and outgrowth \emph{in vitro}: effects of the Arg-Gly-Asp sequence. \emph{Developmental Biology} (1988).
\bibitem{ewald2015} Ewald CY, et al. Dauer-independent insulin/IGF-1-signalling implicates collagen remodelling in longevity. \emph{Nature} (2015).
\bibitem{hansen2014} Hansen M, Kjaer M. Influence of sex and estrogen on musculotendinous protein turnover at rest and after exercise. \emph{Exercise and Sport Sciences Reviews} (2014). \emph{[Added in expansion; verify.]}
\bibitem{heinemeier2013} Heinemeier KM, et al. Lack of tissue renewal in human adult Achilles tendon is revealed by nuclear bomb $^{14}$C. \emph{FASEB Journal} (2013). \emph{[Added in expansion; verify.]}
\bibitem{jamon2014} Jamon M. The development of vestibular system and related functions in mammals: impact of gravity. \emph{Frontiers in Integrative Neuroscience} (2014).
\bibitem{legoff2006} Le Goff C, et al. Regulation of procollagen amino-propeptide processing during mouse embryogenesis by specialization of homologous ADAMTS proteases. \emph{Development} (2006).
\bibitem{maquart2005} Maquart FX, et al. Matrikines in the regulation of extracellular matrix degradation. \emph{Biochimie} (2005).
\bibitem{nowlan2010} Nowlan NC. Mechanobiology of embryonic skeletal development. \emph{Birth Defects Research C} (2010).
\bibitem{pagnon2008} Pagnon-Minot A, et al. Collagen XV, a novel factor in zebrafish notochord differentiation and muscle development. \emph{Developmental Biology} (2008).
\bibitem{panebianco2025} Panebianco CJ, et al. Maternal exercise rescues fetal akinesia-impaired joint and bone morphogenesis. \emph{FASEB Journal} (2025).
\bibitem{peng2022} Peng Q, et al. The sex determination gene \emph{doublesex} regulates expression and secretion of the basement membrane protein Collagen IV. \emph{Journal of Genetics and Genomics} (2022).
\bibitem{sell2005} Sell DR, et al. Glucosepane is a major protein cross-link of the senescent human extracellular matrix. \emph{Journal of Biological Chemistry} (2005). \emph{[Added in expansion; verify.]}
\bibitem{sherwood2004} Sherwood OD. Relaxin's physiological roles and other diverse actions. \emph{Endocrine Reviews} (2004). \emph{[Added in expansion; verify.]}
\bibitem{teuscher2024} Teuscher AC, et al. Longevity interventions modulate mechanotransduction and extracellular matrix homeostasis in \emph{C.\ elegans}. \emph{Nature Communications} (2024).
\bibitem{verbruggen2016} Verbruggen SW, et al. Modeling the biomechanics of fetal movements. \emph{Biomechanics and Modeling in Mechanobiology} (2016).
\bibitem{verzijl2000} Verzijl N, et al. Effect of collagen turnover on the accumulation of advanced glycation end products. \emph{Journal of Biological Chemistry} (2000). \emph{[Added in expansion; verify.]}
\bibitem{word2007} Word RA, et al. Dynamics of cervical remodeling during pregnancy and parturition: mechanisms and current concepts. \emph{Seminars in Reproductive Medicine} (2007). \emph{[Added in expansion; verify.]}
\bibitem{zhou2020} Zhou T, et al. Piezo1/2 mediate mechanotransduction essential for bone formation through concerted activation of NFAT-YAP1-$\beta$-catenin. \emph{eLife} (2020).
\end{thebibliography}

\end{document}
